% ============================================================================
% fig:gapscaling -- finite-size scaling of the two gaps of the 1D Hubbard ring
% AUTO-GENERATED by make_fig_gapscaling.py.  DO NOT EDIT BY HAND.
% Every number below is read at build time from:
%   gap_scaling.json   (gap_scaling.py, 2026-09-18T12:01:30Z)
%   data/spinqw_L12.json      -> eta = 0.10  (spin channel, S^zz(q,w))
%   data/akw_lanczos_L12.json -> eta = 0.18  (one-particle channel, A(k,w))
% Engine: akw_lanczos; full reorthogonalization (2 DGKS passes); broadening: none -- exact poles from the Lanczos tridiagonal.
% Palette inherited from fig_scaling2_native.tex (scMag / scFrac / inkPrim / scDim).
% ---------------------------------------------------------------------------
% PROVENANCE, quantity by quantity (all paths inside gap_scaling.json unless noted):
%   (a) spin gaps          -> /hubbard[L]/spin_gap
%   (a) Heisenberg control -> /heisenberg_control[L]/szz_pi_lowest_pole
%   (a) dashed fit + open square intercept -> /fits/spin_L6plus/{c0_intercept,c0_se,c1_slope}
%   (a) grey band half-width eta           -> release/data/spinqw_L12.json /eta
%   (b) charge gaps        -> /hubbard[L]/charge_gap  ( = mu_plus - mu_minus, re-checked)
%   (b) dashed  curve      -> /fits/charge_L6plus/{c0_intercept,c1_slope}          (L >= 6)
%   (b) dash-dot-dot curve -> /fits/charge_L6plus/quadratic/{c0_intercept,c1,c2}   (L >= 6)
%   (b) dash-dot curve     -> /fits/charge_all/quadratic/{c0_intercept,c1,c2}      (ALL five L)
%   (b) dotted Lieb-Wu line-> /lieb_wu/form_A_sinh_integral
%   (c) L*Delta            -> /hubbard[L]/{L_times_spin_gap,L_times_charge_gap}
%   (c) thin straight lines and quoted slopes -> /L_times_gap/{spin,charge}/{d0,d1,d1_se}
%   (d) the three charge intercepts and their standard errors ->
%         /fits/charge_L6plus/{c0_intercept,c0_se},
%         /fits/charge_L6plus/quadratic/{c0_intercept,c0_se},
%         /fits/charge_all/quadratic/{c0_intercept,c0_se}; dotted line -> /lieb_wu.
% NOTE ON FIT PROTOCOL: every fit drawn here excludes L=4 EXCEPT the dash-dot curve, which is
% the only all-sizes fit and is labelled as such in the legend, in (d) and in the caption.
% Axis limits, tick positions and label positions are layout choices, not data.
% ============================================================================
\definecolor{scMag}{HTML}{D55E00}\definecolor{scFrac}{HTML}{0072B2}
\definecolor{inkPrim}{HTML}{1B1B1B}\definecolor{scDim}{HTML}{8A8A8A}
\definecolor{scCtrl}{RGB}{123,63,160}
\begin{tikzpicture}
\begin{groupplot}[group style={group size=2 by 2, horizontal sep=1.55cm, vertical sep=2.05cm},
  width=7.00cm, height=5.35cm, paperaxis,
  title style={at={(0,1)},anchor=south west,font=\small\bfseries,yshift=1.5pt},
  label style={font=\small}, tick label style={font=\footnotesize},
  xticklabel style={/pgf/number format/fixed}]
% ---------- (a) the spin channel ----------
\nextgroupplot[title={(a)~spin channel},
  xlabel={$1/L$}, ylabel={$\Delta_{\mathrm{s}}\;[t]$},
  xmin=-0.014, xmax=0.278, ymin=-0.062, ymax=0.655,
  xtick={0,0.05,0.10,0.15,0.20,0.25},
  ytick={0,0.1,0.2,0.3,0.4,0.5}, ymajorgrids, grid style={draw=black!10}]
\fill[scDim,fill opacity=0.18] (axis cs:-0.014,0) rectangle (axis cs:0.278,0.2000);
\draw[black!45,line width=0.4pt] (axis cs:-0.014,0) -- (axis cs:0.278,0);
\node[scDim!60!black,font=\tiny,anchor=east,align=right] at (axis cs:0.271,0.108)
  {linewidth\\$2\eta=0.20\,t$};
\node[scFrac,font=\tiny,anchor=west] at (axis cs:0.002,0.612) {Hubbard $U/t=8$};
\node[scCtrl,font=\tiny,anchor=west] at (axis cs:0.002,0.540) {Heisenberg $J=4t^{2}/U$};
\addplot[scFrac,line width=0.9pt,dashed] coordinates {(0,-0.019610) (0.166667,0.342669)};
\addplot[scFrac,only marks,mark=*,mark size=2.4pt,
  mark options={fill=scFrac,draw=white,line width=0.6pt}] coordinates {(0.166667,0.348566) (0.125000,0.239584) (0.100000,0.199559) (0.083333,0.166346)};
\addplot[scFrac,only marks,mark=o,mark size=2.4pt,line width=0.8pt] coordinates {(0.250000,0.332317)};
\addplot[scFrac,only marks,mark=square*,mark size=2.1pt,
  mark options={fill=white,draw=scFrac,line width=0.8pt},
  error bars/.cd,y dir=both,y explicit,error bar style={scFrac,line width=0.7pt},
  error mark options={rotate=90,scFrac,mark size=1.6pt,line width=0.7pt}]
  coordinates {(0,-0.019610) +- (0,0.020421)};
% control drawn last, as OPEN triangles: at L=6 the two channels coincide to 0.006 t
\addplot[scCtrl,only marks,mark=triangle,mark size=3.4pt,line width=0.8pt]
  coordinates {(0.250000,0.500000) (0.166667,0.342371) (0.125000,0.261337) (0.100000,0.211620) (0.083333,0.177924)};
% ---------- (b) the charge channel ----------
\nextgroupplot[title={(b)~charge channel},
  xlabel={$1/L$}, ylabel={$\Delta_{\mathrm{c}}=\mu^{+}-\mu^{-}\;[t]$},
  xmin=-0.014, xmax=0.278, ymin=4.30, ymax=6.34,
  xtick={0,0.05,0.10,0.15,0.20,0.25},
  ytick={4.5,5.0,5.5,6.0}, ymajorgrids, grid style={draw=black!10},
  legend cell align=left, legend columns=1,
  legend style={at={(0.025,0.975)},anchor=north west,font=\tiny,
    draw=black!25,fill=white,fill opacity=0.95,text opacity=1,
    inner sep=1.6pt,row sep=0.4pt}]
\addplot[scMag,line width=0.9pt,dashed] coordinates {(0,4.560890) (0.166667,5.364891)};
\addlegendentry{linear, $L\ge6$}
\addplot[scMag!60,line width=0.9pt,dash dot dot] coordinates {(0.000000,4.461071) (0.004167,4.488103) (0.008333,4.514902) (0.012500,4.541469) (0.016667,4.567804) (0.020833,4.593905) (0.025000,4.619775) (0.029167,4.645411) (0.033333,4.670815) (0.037500,4.695987) (0.041667,4.720926) (0.045833,4.745632) (0.050000,4.770105) (0.054167,4.794347) (0.058333,4.818355) (0.062500,4.842131) (0.066667,4.865674) (0.070833,4.888985) (0.075000,4.912063) (0.079167,4.934909) (0.083333,4.957522) (0.087500,4.979902) (0.091667,5.002050) (0.095833,5.023965) (0.100000,5.045647) (0.104167,5.067097) (0.108333,5.088315) (0.112500,5.109300) (0.116667,5.130052) (0.120833,5.150571) (0.125000,5.170858) (0.129167,5.190913) (0.133333,5.210735) (0.137500,5.230324) (0.141667,5.249681) (0.145833,5.268805) (0.150000,5.287696) (0.154167,5.306355) (0.158333,5.324781) (0.162500,5.342975) (0.166667,5.360936)};
\addlegendentry{$+\,1/L^{2}$, $L\ge6$}
\addplot[scMag,line width=0.9pt,dash dot] coordinates {(0.000000,4.764143) (0.004167,4.769742) (0.008333,4.775843) (0.012500,4.782445) (0.016667,4.789549) (0.020833,4.797155) (0.025000,4.805262) (0.029167,4.813871) (0.033333,4.822982) (0.037500,4.832595) (0.041667,4.842709) (0.045833,4.853325) (0.050000,4.864442) (0.054167,4.876062) (0.058333,4.888183) (0.062500,4.900806) (0.066667,4.913930) (0.070833,4.927557) (0.075000,4.941685) (0.079167,4.956314) (0.083333,4.971446) (0.087500,4.987079) (0.091667,5.003213) (0.095833,5.019850) (0.100000,5.036988) (0.104167,5.054628) (0.108333,5.072770) (0.112500,5.091413) (0.116667,5.110558) (0.120833,5.130205) (0.125000,5.150353) (0.129167,5.171004) (0.133333,5.192156) (0.137500,5.213809) (0.141667,5.235964) (0.145833,5.258621) (0.150000,5.281780) (0.154167,5.305441) (0.158333,5.329603) (0.162500,5.354267) (0.166667,5.379432) (0.170833,5.405100) (0.175000,5.431269) (0.179167,5.457939) (0.183333,5.485112) (0.187500,5.512786) (0.191667,5.540962) (0.195833,5.569639) (0.200000,5.598819) (0.204167,5.628500) (0.208333,5.658682) (0.212500,5.689367) (0.216667,5.720553) (0.220833,5.752241) (0.225000,5.784430) (0.229167,5.817121) (0.233333,5.850314) (0.237500,5.884009) (0.241667,5.918206) (0.245833,5.952904) (0.250000,5.988103)};
\addlegendentry{$+\,1/L^{2}$, all $L$}
\addplot[inkPrim,line width=0.7pt,densely dotted] coordinates {(-0.014,4.679517) (0.278,4.679517)};
\addlegendentry{Lieb--Wu $\Delta_{\infty}$}
\addplot[forget plot,scMag,only marks,mark=*,mark size=2.4pt,
  mark options={fill=scMag,draw=white,line width=0.6pt}] coordinates {(0.166667,5.358127) (0.125000,5.185842) (0.100000,5.022235) (0.083333,4.968759)};
\addplot[forget plot,scMag,only marks,mark=o,mark size=2.4pt,line width=0.8pt] coordinates {(0.250000,5.991359)};
% ---------- (c) opposite finite-size behaviour, model free ----------
% the two thin straight lines are the SAME linear fits whose slopes are quoted; they are
% sampled at 41 points each because a straight line in L is curved on a logarithmic y axis.
\nextgroupplot[title={(c)~opposite size dependence of $L\Delta$},
  xlabel={$L$}, ylabel={$L\,\Delta\;[t]$}, ymode=log,
  xmin=3.2, xmax=13.0, ymin=0.92, ymax=420,
  xtick={4,6,8,10,12}, ytick={1,2,5,10,20,50,100,200},
  yticklabels={1,2,5,10,20,50,100,200},
  ymajorgrids, grid style={draw=black!10}]
\addplot[scMag!55,line width=0.5pt,forget plot] coordinates {(6.0000,32.196040) (6.1500,32.879775) (6.3000,33.563510) (6.4500,34.247245) (6.6000,34.930980) (6.7500,35.614715) (6.9000,36.298450) (7.0500,36.982185) (7.2000,37.665921) (7.3500,38.349656) (7.5000,39.033391) (7.6500,39.717126) (7.8000,40.400861) (7.9500,41.084596) (8.1000,41.768331) (8.2500,42.452066) (8.4000,43.135801) (8.5500,43.819536) (8.7000,44.503271) (8.8500,45.187006) (9.0000,45.870741) (9.1500,46.554476) (9.3000,47.238211) (9.4500,47.921946) (9.6000,48.605681) (9.7500,49.289416) (9.9000,49.973151) (10.0500,50.656886) (10.2000,51.340621) (10.3500,52.024356) (10.5000,52.708091) (10.6500,53.391826) (10.8000,54.075561) (10.9500,54.759296) (11.1000,55.443031) (11.2500,56.126766) (11.4000,56.810501) (11.5500,57.494236) (11.7000,58.177971) (11.8500,58.861706) (12.0000,59.545441)};
\addplot[scFrac!55,line width=0.5pt,forget plot] coordinates {(6.0000,2.030974) (6.1500,2.029423) (6.3000,2.027872) (6.4500,2.026320) (6.6000,2.024769) (6.7500,2.023218) (6.9000,2.021667) (7.0500,2.020116) (7.2000,2.018565) (7.3500,2.017014) (7.5000,2.015463) (7.6500,2.013912) (7.8000,2.012361) (7.9500,2.010810) (8.1000,2.009259) (8.2500,2.007708) (8.4000,2.006157) (8.5500,2.004606) (8.7000,2.003055) (8.8500,2.001504) (9.0000,1.999953) (9.1500,1.998402) (9.3000,1.996850) (9.4500,1.995299) (9.6000,1.993748) (9.7500,1.992197) (9.9000,1.990646) (10.0500,1.989095) (10.2000,1.987544) (10.3500,1.985993) (10.5000,1.984442) (10.6500,1.982891) (10.8000,1.981340) (10.9500,1.979789) (11.1000,1.978238) (11.2500,1.976687) (11.4000,1.975136) (11.5500,1.973585) (11.7000,1.972034) (11.8500,1.970483) (12.0000,1.968931)};
\addplot[scMag,only marks,mark=*,mark size=2.4pt,
  mark options={fill=scMag,draw=white,line width=0.6pt}] coordinates {(6.000000,32.148761) (8.000000,41.486737) (10.000000,50.222353) (12.000000,59.625112)};
\addplot[scMag,only marks,mark=o,mark size=2.4pt,line width=0.8pt] coordinates {(4.000000,23.965437)};
\addplot[scFrac,only marks,mark=*,mark size=2.4pt,
  mark options={fill=scFrac,draw=white,line width=0.6pt}] coordinates {(6.000000,2.091398) (8.000000,1.916668) (10.000000,1.995590) (12.000000,1.996155)};
\addplot[scFrac,only marks,mark=o,mark size=2.4pt,line width=0.8pt] coordinates {(4.000000,1.329266)};
\node[scMag,font=\tiny,anchor=west,align=left] at (axis cs:3.42,235)
  {charge: fitted slope in $L$, $+4.56\pm0.05$};
\node[scFrac,font=\tiny,anchor=west,align=left] at (axis cs:3.42,6.7)
  {spin: fitted slope in $L$, $-0.010\pm0.018$};
% ---------- (d) the three L->inf charge intercepts, with their standard errors ----------
% Same three fits as (b); this panel exists because on the scale of (b) a +-0.044 t error bar
% is thinner than its own marker, so the 2.70 s.e. tension with Lieb-Wu cannot be checked by eye.
\nextgroupplot[title={(d)~the $1/L\to0$ charge intercepts},
  xlabel={$\Delta_{\mathrm{c}}(L\to\infty)\;[t]$}, ymin=0.15, ymax=3.75,
  xmin=3.96, xmax=4.99, xtick={4.0,4.2,4.4,4.6,4.8}, ytick=\empty,
  xmajorgrids, grid style={draw=black!10}]
\addplot[inkPrim,line width=0.7pt,densely dotted,forget plot] coordinates {(4.679517,0.75) (4.679517,3.75)};
\node[inkPrim,font=\tiny,anchor=north] at (axis cs:4.679517,0.62) {Lieb--Wu $4.6795\,t$};
\addplot[scMag,only marks,mark=square*,mark size=2.3pt,
  mark options={fill=white,draw=scMag,line width=0.8pt},
  error bars/.cd,x dir=both,x explicit,error bar style={scMag,line width=0.7pt},
  error mark options={scMag,mark size=2.0pt,line width=0.7pt}]
  coordinates {(4.560890,3.00) +- (0.043868,0)};
\node[scMag,font=\tiny,anchor=south west,align=left] at (axis cs:3.985,3.33)
  {linear in $1/L$, $L\ge6$};
\node[scDim!55!black,font=\tiny,anchor=south west,align=left] at (axis cs:3.985,3.11)
  {$2.70$ s.e.\ below Lieb--Wu};
\addplot[scMag,only marks,mark=triangle*,mark size=3.0pt,
  mark options={fill=white,draw=scMag,line width=0.8pt},
  error bars/.cd,x dir=both,x explicit,error bar style={scMag,line width=0.7pt},
  error mark options={scMag,mark size=2.0pt,line width=0.7pt}]
  coordinates {(4.461071,2.00) +- (0.308445,0)};
\node[scMag,font=\tiny,anchor=south west,align=left] at (axis cs:3.985,2.33)
  {$+\,1/L^{2}$, $L\ge6$};
\node[scDim!55!black,font=\tiny,anchor=south west,align=left] at (axis cs:3.985,2.11)
  {$0.71$ s.e.\ below Lieb--Wu};
\addplot[scMag,only marks,mark=diamond*,mark size=3.0pt,
  mark options={fill=white,draw=scMag,line width=0.8pt},
  error bars/.cd,x dir=both,x explicit,error bar style={scMag,line width=0.7pt},
  error mark options={scMag,mark size=2.0pt,line width=0.7pt}]
  coordinates {(4.764143,1.00) +- (0.127814,0)};
\node[scMag,font=\tiny,anchor=south west,align=left] at (axis cs:3.985,1.33)
  {$+\,1/L^{2}$, all five $L$\\(includes $L=4$)};
\node[scDim!55!black,font=\tiny,anchor=south west,align=left] at (axis cs:3.985,1.11)
  {$0.66$ s.e.\ above Lieb--Wu};
\end{groupplot}
\end{tikzpicture}
